\statement{B.79}{Proved}{The actual marginal has a source-marked presentation with no retained base loops}
The same actual marginal group and source of B.78 admit a presentation
on twenty-three fiber generators with eighty-two local relation normal
generators. The base of this pencil is \(\mathbb C\), so every monodromy
relation is included; there is no further retained-base saturation to
perform. The entire source deletion-kernel equality is equivalent to
normalization by three specified retained directions after one local
elimination. These statements hold for actual compact images in full
completions. The normalization condition itself remains unresolved.

\begin{proof}
\emph{The complete sum pencil.}\par
In the normalized ordered surface of B.78 put
\[
              T=\frac{v+w}{2},\qquad z=w-1.
 \tag{B.79.1}
\]
The generic fiber has twenty-three finite punctures:
\[
\begin{split}
 &-1,\quad 2T-1,\quad T-1;\\
 &u-1,\quad 2T-1-u\qquad(u\in\mu _5);\\
 &z=(T-1)\pm\sqrt{T^2-q}\qquad(q\in\mu _5).
\end{split}
 \tag{B.79.2}
\]
The eight deleted punctures are the constant and moving terms with
\(u\ne1\). Every component has finite degree over \(\mathbb C\),
and no entire fiber is deleted. There are exactly twenty-six
exceptional values. Their local types and relation counts are:
\[
\begin{array}{c|c|c|c|c}
\text{value}&\text{number}&\text{local braid}&
\text{exponent}&\text{relations each}\\\hline
0&1&\Delta_3^2&6&2\\
u\in\mu _5&5&\Delta_5\Delta_{3,\mathrm{inner}}&13&4\\
-u,\ u\in\mu _5&5&\Delta_3&3&2\\
u/2&5&\Delta_2^2\,\Delta_2^2&4&2\\
(u+v)/2,\ u\ne v&10&\Delta_3^2\,\Delta_3^2&12&4
\end{array}
 \tag{B.79.3}
\]
Products in the last two rows have disjoint supports. The symbols
\(\Delta_n\) denote positive half twists on the indicated consecutive
local strands. The fiber losses in the five rows are respectively
\(2,4,2,2,4\); their sum is \(82\), giving
\(\chi(E)=1-23+82=60\).

Here is an exact completeness check, not just a numerical count.
The thirteen rational punctures give linear pair differences. A
quadratic evaluated at a rational puncture has degree at most two in
\(T\), and its own discriminant is \(4(T^2-q)\). Two different
quadratics have constant nonzero difference \(q-q'\), so their roots
never collide. For every remaining collision polynomial, the exact
sum of its orders at the values in (B.79.3) equals its full degree.
Arithmetic in \(\mathbb Q[\zeta]\) verifies that all twenty-six
values are distinct. The discriminant exponent is twice the sum of
the rational-pair and rational/quadratic resultant orders, plus the
quadratic discriminant orders. This gives the fourth column of
(B.79.3) independently of the braid extraction. The complete calculation
is \path{work_DE/next_marginal_source/audit_sum_pencil.py}, with genuine
output and \path{work_DE/next_marginal_source/sum_pencil_critical_audit.json}.

\emph{One based marking, including the special source.}\par
Take \(T_*=1+i/100,z_*=100\) and projection \(Pz\) with the same
\(P=1+371i/1000\) as B.78. Use exactly the original source path and
root ordering of B.78, then follow \(T=1+iy,z=100\),
\(0\leq y\leq1/100\), lifting \(s^2=vw\). The source is \(T=1\).
Along this short path a fixed disk \(\{|Pz|<1/4\}\) isolates its
five merging retained points. The other eighteen punctures stay outside
the disk. Based cuts outside it and a common boundary cut therefore
give the literal source boundary
\[
               C=y_{10}y_{11}y_{12}y_{13}y_{14}.
 \tag{B.79.4}
\]
The initial strand-label order, using rational labels zero through
twelve and quadratic labels thirteen through twenty-two in the order
of (B.79.2), is
\[
 (5,6,18,20,16,4,0,22,7,14,8,2,3,13,12,21,1,9,15,19,17,11,10).
 \tag{B.79.5}
\]
Thus the five cluster positions are exactly \(10,\ldots,14\).
Replace \(y_{10}\) by \(x_{10}=C\), leaving \(x_i=y_i\) for
\(i\ne10\). Both substitutions are explicitly checked. The actual
source is the free factor
\[
 H=F(x_i:i\notin\{11,12,13,14\})\cong F_{19},
 \qquad \mathcal U=\{1,2,6,9,15,18,22,23\}.
 \tag{B.79.6}
\]
Here \(\mathcal U\) is the set of eight deleted source meridians.
The two finite-cover kernel comparisons in B.78 identify this source
and its actual kernel image with the original marginal, including in
full completions.

All twenty-six lassos use disjoint squares of half-side \(1/1000\)
and straight whiskers from \(T_*\) to their northeast corners. Exact
distance and determinant checks prove that these are a distinguished
system. The rational coefficient deformation, Sturm isolation and
height bounds of B.78 apply to this quadratic family as well. One
additional projection issue is essential: the diagonal point is the
midpoint of a quadratic pair, so a pair's vertical crossing forces
a three-strand vertical alignment for every linear projection.

The new engine certifies distinct imaginary heights for every pair
in each connected inversion cluster. In a sufficiently small interval
around its isolated event, the strands remain in disjoint height
disks. The product of those disks is simply connected. Replacing the
horizontal projection there by an adjacent-inversion sorting path,
with the same certified height order, is consequently an isotopy
relative to endpoints. Every inverted pair crosses once; its sign is
fixed by its height difference. This resolves the nongeneric projection
without changing the geometric braid. The initial strict-run failure
and the successful rerun are both retained.

The full words and all 130 segment certificates are in
\path{work_DE/next_marginal_source/sum_pencil_braids.json}; the producer
is \path{work_DE/next_marginal_source/sum_pencil_braid.py}, with genuine
\texttt{.out} and \texttt{.err} files. The exact source-cluster audit is
\path{work_DE/next_marginal_source/check_sum_source_cluster.py}.
In particular, both the actual local and global source lasso have
Artin action \(\Delta_5\Delta_{3,\mathrm{inner}}\), and their common
whisker fixes the word \(C\) in (B.79.4).

\emph{The eighty-two actual relation generators.}\par
Use the pullback convention of B.78. For each lasso the program absorbs
an explicitly recorded initial local braid prefix into its whisker
\(W\), obtaining precisely one supported standard form in (B.79.3).
It checks the equality of the full global Artin action with
\(W\beta_{\mathrm{standard}}W^{-1}\), as well as both inverse
whisker actions. Three negative-root cases require a two-letter
prefix; the other twenty-three cases require none.

For a full twist on local meridians \(m_1,\ldots,m_n\), write
\(P=m_1\cdots m_n\) and use the \(n-1\) relations
\([P,m_j]=1\), \(j<n\). For a three-point half twist use
\(m_3=m_1\) and \([m_1,m_2]=1\). For the five-point form put
\(R=m_1m_2m_3m_4\) and use
\[
       m_5=m_1,\qquad [R,m_2]=[R,m_3]=[R,m_4]=1.
 \tag{B.79.7}
\]
Indeed its Artin action sends \(m_1\) to \(Rm_5R^{-1}\),
\(m_5\) to \(m_1\), and the three inner meridians to their
\(R\)-conjugates. Relations (B.79.7) also imply \([R,m_1]=1\),
so they are equivalent to all its variation relations. These formulas,
transported by \(\Phi_W\) and rewritten in the source basis, give
\(2+20+10+10+40=82\) relation normal generators.

The complete list is
\path{work_DE/next_marginal_source/sum_pencil_presentation.json}, produced
by \path{work_DE/next_marginal_source/extract_sum_presentation.py} and
its genuine output. It specifies
\[
             G\cong F(x_1,\ldots,x_{23})/
                         \langle\!\langle\mathscr R_{82}\rangle\!\rangle.
 \tag{B.79.8}
\]
A far-point section exists on all \(\mathbb C\). Ordinary van Kampen
gluing, or the same generic-fiber theorem checked in B.78, therefore
imposes all twenty-six variation families and leaves no base generator.
The eighty-two count refers to locally minimal normal generating
sets for the displayed variation families. Neither the twenty-three
ambient generators nor the eighty-two global relations are claimed
to be globally minimal. Completing (B.79.8) gives the corresponding quotient
of the free profinite group by the closed normal closure of these
same relations.

\emph{The whole kernel condition and three retained directions.}\par
Define \(N=\langle\!\langle x_i:i\in\mathcal U\rangle\!\rangle_H\),
let \(I\) be its actual image in \(G\), and let \(R_U\) be the
ambient normal closure of the eight deleted meridians. Their actual
completed images are \(\mathcal I=\overline I\) and
\(\mathcal R=\overline{R_U}\), as in B.78. Since \(H\) and its four
missing retained generators generate \(G\), equality \(I=R_U\)
is equivalent to invariance of \(I\) under all four of those generators;
the same equivalence holds for the closed images.

One local elimination reduces this to three directions. In the actual
source block put \(a=y_{10}, b=y_{11}, d=y_{12}\). Its variation
relations give
\[
 \begin{split}
 y_{14}&=a,\qquad
 y_{13}=d^{-1}b^{-1}a^{-1}Ca^{-1},\\
 [C,a]&=1,\qquad [Ca^{-1},b]=[Ca^{-1},d]=1.
 \end{split}
 \tag{B.79.9}
\]
Thus the quotient defined by this local family alone is
\(\mathbb Z(R)\times F_3(a,b,d)\), with
\(R=Ca^{-1}\) central and \(C=Ra\). The source meridian is \(C\),
not \(R\). No injection of this local quotient into the full ambient
group is inferred. The ambient group is generated by the actual source image
and \(a,b,d\), so the precise remaining condition is
\[
 \mathcal I=\mathcal R
 \quad\Longleftrightarrow\quad
 a\mathcal I a^{-1}=b\mathcal I b^{-1}
                       =d\mathcal I d^{-1}=\mathcal I.
 \tag{B.79.10}
\]
The discrete version is identical. For compact images, closedness of
the normalizer supplies all limits of words in these generators.
Neither internal normal generation of \(\mathcal R\) by the source
image nor checks in nilpotent quotients establish (B.79.10). The whole
normalization condition, rather than just the designated peripheral
generators or their first few conjugates, remains the main marginal
source-membership problem.
\end{proof}
